\documentclass[12pt,a4paper]{article}
\usepackage[utf8]{inputenc}
\usepackage[T1]{fontenc}
\usepackage[spanish]{babel}
\usepackage{amsmath, amssymb}
\usepackage{graphicx}
\usepackage{booktabs}
\usepackage{float}
\usepackage{hyperref}
\usepackage{geometry}
\geometry{margin=2.5cm}

\begin{document}

\begin{titlepage}
    \centering
    \vspace*{2cm}
    {\Large Cálculo Aplicado\par}
    \vspace{2cm}
    {\LARGE \textbf{Optimización Numérica en una dimensión}\par}
    \vspace{2cm}
    {\large 11 de junio de 2026\par}
    \vspace{2cm}
    {\large \textbf{Estudiantes:}\par}
    \vspace{0.5cm}
    {\large Mateo Hernández (5.565.681-8)\par}
    {\large Agustín Pose (5.386.051-8)\par}
    \vfill
\end{titlepage}

\newpage

\tableofcontents
\newpage

\begin{abstract}
Este informe presenta la implementación y análisis de tres métodos clásicos de optimización numérica en una dimensión: bisección, Newton y descenso por gradiente. Se verifica su funcionamiento en funciones de prueba simples y en un problema de comparación de complejidad computacional. Además, se aplican a problemas de ajuste de datos con tres modelos distintos. Se incluyen extensiones (SGD, learning rate variable, formulación producto) que ilustran técnicas modernas de optimización.
\end{abstract}
\newpage

\section{Introducción}
La optimización numérica es fundamental para resolver problemas donde no existe solución analítica. En una dimensión, permite aproximar puntos de mínimo o máximo de forma iterativa. Este trabajo estudia tres métodos clásicos y sus variantes modernas.

\textbf{Objetivos:}
\begin{enumerate}
    \item Implementar bisección, Newton y descenso por gradiente
    \item Verificar su comportamiento en funciones de prueba
    \item Analizar la comparación de complejidad: $f(n)=n^2-50n\log_2(n)$
    \item Resolver problemas de ajuste de datos con tres modelos
    \item Explorar extensiones: SGD, learning rate variable, formulación producto
\end{enumerate}

\section{Marco teórico}
Un punto crítico de $f(x)$ satisface $f'(x)=0$. Su clasificación usa el signo de $f''(x)$.

\textbf{Bisección:} Busca cambio de signo en $f'(x)$ dentro de $[a,b]$ y reduce el intervalo iterativamente.

\textbf{Newton:} Usa la aproximación local $f(x) \approx f(x_n) + f'(x_n)(x-x_n) + \frac{1}{2}f''(x_n)(x-x_n)^2$. Al igualar la derivada a cero:
$$x_{n+1} = x_n - \frac{f'(x_n)}{f''(x_n)}$$

\textbf{Descenso por gradiente:} Actualiza siguiendo $x_{n+1} = x_n - \alpha f'(x_n)$ donde $\alpha$ es el learning rate.

\textbf{SGD:} Versión estocástica del descenso por gradiente, útil en minibatches.

\textbf{Learning rate variable:} Usa schedule $\alpha_n = \alpha_0 \cdot \text{decay}(n)$ para mejorar convergencia.

\textbf{Formulación producto:} Maximiza $P(w) = \prod \exp(-(y_i - \hat{y}_i)^2)$ equivalente a minimizar la suma de cuadrados en el dominio log.

\section{Metodología}
\subsection{Implementación}
Se utilizó un framework simbólico que permite:
\begin{itemize}
    \item Definir funciones mediante sobrecarga de operadores: $f = x^2 + 2\sin(3x)$
    \item Calcular derivadas analíticamente: $f' = 2x + 6-\cos(3x)$, $f'' = 2 - 18\sin(3x)$
    \item Evaluar en cualquier punto sin necesidad de derivadas numéricas
\end{itemize}

\subsection{Estrategia de verificación}
\begin{enumerate}
    \item Funciones simples con extremos conocidos: $f(x)=x^2$, $f(x)=(x+1/2)^3-(x+1/2)$, $f(x)=-\cos(x)$
    \item Función computacional: $f(n)=n^2-50n\log_2(n)$ (complejidad de algoritmos)
    \item Datos sintéticos con tres modelos: lineal $y=2x$, exponencial $y=e^{x/4.5}$, seno $y=\sin(1.5x)$
\end{enumerate}


\paragraph{Repositorio del proyecto}
\url{https://github.com/AgustinPose/Optimizaci-n-Num-rica-en-una-Dimensi-n}

\section{Resultados}
\subsection{Verificación en funciones simples}
Se probaron los tres métodos en funciones elementales:

\begin{center}
\begin{tabular}{lrrrr}
\toprule
Función & $f(x^*)$ & Bisección & Newton & GD \\
\midrule
$x^2$ & $0$ en $x=0$ & \checkmark & \checkmark & \checkmark \\
$(x+1/2)^3-(x+1/2)$ & $0$ en $x=-1/2$ & \checkmark & \checkmark & \checkmark \\
$-\cos(x)$ & $0$ en $x=\pi/2$ & \checkmark & \checkmark & \checkmark \\
\bottomrule
\end{tabular}
\end{center}

Todos los métodos convergieron en funciones bien comportadas. Newton fue el más rápido (típicamente 3-5 iteraciones). 

\textbf{BONUS: Visualización con historial de convergencia}

Se implementaron versiones de los métodos que devuelven la secuencia de aproximaciones durante la optimización. La visualización muestra el camino recorrido por cada algoritmo con código de colores (azul $\to$ primeras iteraciones, amarillo $\to$ últimas). Este análisis permite observar la velocidad de convergencia de cada método y cómo la trayectoria varía según el algoritmo y los parámetros iniciales.

\subsection{Análisis computacional: $f(n)=n^2-50n\log_2(n)$}
Esta función surge al comparar dos algoritmos con complejidades $T_A(n)=n^2$ y $T_B(n)=50n\log_2(n)$. El objetivo es encontrar el tamaño de entrada donde ambos tienen tiempo comparable.

\textbf{Contexto:} Se busca resolver $f(n)=0$, es decir, encontrar el valor $n$ donde:
$$n^2 = 50n\log_2(n)$$

\textbf{Comparación de convergencia en intervalo [400, 500]:}
\begin{center}
\begin{tabular}{lrrrr}
\toprule
Método & $n^*$ & $f(n^*)$ & Iteraciones & Observaciones \\
\midrule
Bisección & $438.16$ & $\approx 0$ & 9 & Convergencia estable \\
Newton & $438.16$ & $\approx 0$ & 2 & Convergencia rápida \\
GD ($\alpha=0.000005$) & $438.92$ & $\approx 11.83$ & 5 & Convergencia lenta \\
\bottomrule
\end{tabular}
\end{center}

Newton es significativamente más rápido. Bisección es robusto pero requiere intervalo con cambio de signo. Descenso por gradiente es el más lento, requiriendo ajuste cuidadoso del learning rate.

\subsection{Ajuste de datos}
Se ajustaron tres modelos a datos sintéticos:

\subsubsection{Modelo lineal: $\hat{y}=w \cdot x$}
Datos: 50 puntos de $y = 2x + \epsilon$, $\epsilon \sim N(0, 1)$

\textbf{Resultado óptimo:} $w^* = 1.848$ (comparado con $w_{\text{real}} = 2.0$)

Métodos convergieron al mismo óptimo con diferente número de iteraciones.

\subsubsection{Modelo exponencial: $\hat{y}=e^{w \cdot x}$}
Datos: 50 puntos de $y = e^{x/4.5} + \epsilon$

\textbf{Resultado óptimo:} $w^* = 0.222$ (aproximadamente $1/4.5 \approx 0.222$)

La función de pérdida es suave y tiene un mínimo único bien definido.

\subsubsection{Modelo seno: $\hat{y}=\sin(w \cdot x)$}
Datos: 50 puntos de $y = \sin(1.5x) + 0.5\epsilon$

\textbf{Resultado óptimo:} $w^* = 1.502$ (muy cercano a $w_{\text{real}} = 1.5$)

La función de pérdida es no convexa pero presenta mínimos globales claros.

\section{Extensiones (Bonus)}
\subsection{Descenso por gradiente estocástico (SGD)}
Se implementó SGD y se comparó con descenso por gradiente clásico. En este problema, ambos convergieron a los mismos óptimos:
\begin{center}
\begin{tabular}{lrr}
\toprule
Modelo & SGD & GD clásico \\
\midrule
Lineal & $w=1.848$ & $w=1.848$ \\
Seno & $w=0.690$ & $w=0.690$ \\
\bottomrule
\end{tabular}
\end{center}

SGD es más útil cuando el costo de evaluar todos los datos es prohibitivo.

\subsection{Learning rate variable}
Se probaron dos schedules de decaimiento:
\begin{itemize}
    \item Exponencial: $\alpha_n = \alpha_0 \exp(-0.1n)$
    \item Inverso: $\alpha_n = \frac{\alpha_0}{1+0.1n}$
\end{itemize}

Todos convergieron al mismo punto $x^* \approx -0.471$ con convergencia comparable. El learning rate variable mejora la estabilidad en problemas mal condicionados.

\subsection{Formulación producto}
Se maximizó $P(w) = \prod \exp(-(y_i-\hat{y}_i)^2)$ equivalente a minimizar la suma de cuadrados.

En el dominio log, Newton converge en 1 iteración, mostrando la estabilidad numérica del log-dominio. Descenso por gradiente en el dominio log requiere learning rate muy pequeño para evitar divergencia.

\textbf{Conclusión sobre formulación producto:}
\begin{center}
\begin{tabular}{lrr}
\toprule
Formulación & Método & Convergencia \\
\midrule
Suma & Newton & $w=1.848$, 18 iteraciones \\
Suma & GD & $w=1.848$, 18 iteraciones \\
Producto (log) & Newton & $w=1.848$, 1 iteración \\
\bottomrule
\end{tabular}
\end{center}

\section{Discusión}
\subsection{Elección de método}
\begin{itemize}
    \item \textbf{Bisección:} Cuando se garantiza cambio de signo, muy robusto
    \item \textbf{Newton:} Cuando hay buena condición inicial y se pueden evaluar segunda derivada
    \item \textbf{GD:} Cuando el problema es de alta dimensión (escalable)
    \item \textbf{SGD:} Cuando hay muchos datos y evaluación es costosa
    \item \textbf{Log-dominio:} Cuando estabilidad numérica es crítica
\end{itemize}

\subsection{Impacto del learning rate}
El parámetro $\alpha$ es crítico en descenso por gradiente. Valores muy pequeños requieren más iteraciones; valores muy grandes causan divergencia. Schedules de decaimiento pueden mejorar la convergencia.

\subsection{Tradeoff velocidad-robustez}
Newton es el más rápido pero requiere buena condición inicial y capacidad de evaluar derivadas de segundo orden. Bisección es el más lento pero muy robusto. Descenso por gradiente es un punto intermedio útil en altas dimensiones.

\section{Conclusiones}
\begin{enumerate}
    \item Los tres métodos fueron implementados correctamente y verificados en múltiples funciones
    \item Newton es superior en velocidad pero requiere aproximación cuadrática válida
    \item El descenso por gradiente es más escalable y forma la base de métodos modernos
    \item SGD, learning rate variable y formulación producto son extensiones prácticas útiles
    \item El framework simbólico permitió enfoque limpio sin diferenciación numérica
\end{enumerate}

\section*{Bibliografía}
\begin{itemize}
    \item Nocedal, J., \& Wright, S. (2006). Numerical Optimization. Springer.
    \item Boyd, S., \& Vandenberghe, L. (2004). Convex Optimization. Cambridge University Press.
    \item Goodfellow, I., Bengio, Y., \& Courville, A. (2016). Deep Learning. MIT Press.
\end{itemize}

\end{document}
